% chapters/tikz_ch06.tex - 第 6 章全部插图的高精度原生 TikZ 矢量代码

% 1. 图 6.1: xi_0 的锥形邻域
\newcommand{\drawFigSixOne}{%
\begin{tikzpicture}[scale=0.95, >=Stealth, line join=round]
    \draw[->] (-0.45, 0) -- (4.30, 0) node[right] {$\xi_1$};
    \draw[->] (0, -0.45) -- (0, 3.20) node[above] {$\xi_2$};
    % 锥形邻域 (阴影)
    \fill[gray!18] (0, 0) -- (3.95, 2.36) -- (3.95, 0.97) -- cycle;
    \draw[gray!65] (0, 0) -- (3.95, 2.36);
    \draw[gray!65] (0, 0) -- (3.95, 0.97);
    \node[gray!70!black] at (3.05, 2.28) {\footnotesize 锥形邻域};
    % xi_0
    \draw[->, very thick] (0, 0) -- (2.60, 1.12);
    \node[right] at (2.62, 1.14) {$\xi_0$};
    % 小球 B(xi_0, eps)
    \draw[dashed] (1.05, 0.45) circle (0.45);
    \node[left, inner sep=2pt] at (0.58, 0.52) {\footnotesize $\epsilon$-球};
    \fill (0, 0) circle (1.5pt) node[below left, inner sep=1pt] {$0$};
\end{tikzpicture}%
}

% 2. 图 6.2: H(x1)H(x2) 与 H(x1)delta(x2) 的波前集
\newcommand{\drawFigSixTwo}{%
\begin{tikzpicture}[scale=0.90, >=Stealth, line join=round]
    % 左图: WF(H(x1)H(x2))
    \begin{scope}
        \draw[->] (-2.05, 0) -- (2.45, 0) node[right] {$x_1$};
        \draw[->] (0, -2.05) -- (0, 2.45) node[above] {$x_2$};
        % 半轴 {x2 = 0, x1 > 0} 上的余法方向 (竖直)
        \draw[<->, thick] (1.32, -0.50) -- (1.32, 0.50);
        \draw[<->, thick] (1.75, -0.50) -- (1.75, 0.50);
        % 半轴 {x1 = 0, x2 > 0} 上的余法方向 (水平)
        \draw[<->, thick] (-0.50, 1.32) -- (0.50, 1.32);
        \draw[<->, thick] (-0.50, 1.75) -- (0.50, 1.75);
        % 角点处的方向锥
        \fill[gray!25] (0, 0) -- (0.86, 0.10) -- (0.10, 0.86) -- cycle;
        \node[right, inner sep=2pt] at (0.62, 0.62) {\footnotesize $N^*\{0\}$};
        \node[below] at (0, -2.16) {\footnotesize $H(x_1)H(x_2)$};
    \end{scope}
    % 右图: WF(H(x1)delta(x2))
    \begin{scope}[xshift=6.4cm]
        \draw[->] (-2.05, 0) -- (2.45, 0) node[right] {$x_1$};
        \draw[->] (0, -2.05) -- (0, 2.45) node[above] {$x_2$};
        % 半轴 {x2 = 0, x1 > 0} 上的余法方向 (竖直)
        \draw[<->, thick] (1.32, -0.50) -- (1.32, 0.50);
        \draw[<->, thick] (1.75, -0.50) -- (1.75, 0.50);
        % 原点处的全部方向
        \fill[gray!25] (0, 0) circle (0.34);
        \draw[<->, thick] (-0.92, 0) -- (0.92, 0);
        \node[right, inner sep=2pt] at (0.42, 0.66) {\footnotesize $N^*\{0\}$};
        \node[below] at (0, -2.16) {\footnotesize $H(x_1)\delta(x_2)$};
    \end{scope}
\end{tikzpicture}%
}

% 3. 图 6.3: 微局部椭圆正则性估计 (6.56) 的图示
\newcommand{\drawFigSixThree}{%
\begin{tikzpicture}[scale=0.92, >=Stealth, line join=round]
    % 左图: 相空间中的示意
    \begin{scope}
        \draw[thick] (-0.10, 0) -- (4.60, 0);
        \node[below left] at (0, -0.06) {\footnotesize $S^*M$};
        % Ell(A)
        \draw[line width=7pt, gray!30] (1.05, 0) -- (4.35, 0);
        \draw[line width=3.2pt] (1.95, 0) -- (3.10, 0);
        % 标注 (纵向分层, 互不重叠)
        \node[above] at (2.52, 0.66) {\footnotesize $\|GAu\|_{H^{s-m}} \Rightarrow \|Bu\|_{H^s}$};
        \node[above] at (2.52, 0.24) {\footnotesize $\Ell(A) \supset \Ell(GA)$};
        \node[below] at (2.52, -0.20) {\footnotesize $\WF'(B)$};
    \end{scope}
    % 右图: 核的支集
    \begin{scope}[xshift=7.0cm]
        \draw[thick] (0, 0) rectangle (4.0, 3.4);
        \node[below left] at (0, 0) {\footnotesize $M$};
        \node[left, rotate=90] at (-0.10, 1.7) {\footnotesize $M$};
        % 对角线
        \draw[dashed, gray] (0, 0) -- (3.4, 3.4);
        % K x K 深色方块
        \fill[gray!45] (0.55, 0.55) rectangle (1.95, 1.95);
        \node[below] at (1.25, 0.50) {\footnotesize $K$};
        \node[left] at (0.52, 1.25) {\footnotesize $K$};
        % 带状区域 K x M
        \fill[gray!15] (0.55, 0.55) rectangle (1.95, 3.4);
        % K x K~
        \draw[densely dotted, thick] (0.55, 0.55) rectangle (2.55, 3.05);
        \node[right] at (2.58, 3.05) {\footnotesize $\tilde{K}$};
        \node[above] at (1.25, 2.10) {\footnotesize $K \times \tilde{K}$};
    \end{scope}
\end{tikzpicture}%
}

% 4. 图 6.4: 命题 6.41 中 WF(u), WF(v) 与 A 的微局部性质
\newcommand{\drawFigSixFour}{%
\begin{tikzpicture}[scale=0.98, >=Stealth, line join=round]
    \draw[thick] (0, 0) -- (7.4, 0);
    \node[below left] at (0, 0) {\footnotesize $S^*M$};
    % WF'(A) 与 WF'(I - A) 的分界
    \draw[densely dashed] (3.70, -1.25) -- (3.70, 1.25);
    \node[above] at (3.70, 1.28) {\footnotesize 分界};
    % 左侧: WF'(A) 覆盖, WF(u) 在其中
    \draw[line width=7pt, gray!25] (0.55, 0) -- (3.30, 0);
    \node[above] at (1.90, 0.60) {\footnotesize $\WF'(A)$};
    \draw[line width=3.2pt] (1.10, 0) -- (2.05, 0);
    \node[below] at (1.55, -0.62) {\footnotesize $\WF(u)$};
    \draw[->, thick] (1.55, -0.52) -- (1.55, -0.14);
    % 右侧: WF'(I - A) 覆盖, WF(v) 在其中
    \draw[line width=7pt, gray!25] (4.10, 0) -- (6.85, 0);
    \node[above] at (5.50, 0.60) {\footnotesize $\WF'(I - A)$};
    \draw[line width=3.2pt] (5.30, 0) -- (6.25, 0);
    \node[below] at (5.75, -0.62) {\footnotesize $\WF(v)$};
    \draw[->, thick] (5.75, -0.52) -- (5.75, -0.14);
\end{tikzpicture}%
}

% 5. 图 6.5: 命题 6.44 证明中的若干锥 (n = 2, k = 1)
\newcommand{\drawFigSixFive}{%
\begin{tikzpicture}[scale=1.00, >=Stealth, line join=round]
    \draw[->] (-2.05, 0) -- (4.85, 0) node[right] {$\xi$};
    \draw[->] (0, -2.30) -- (0, 2.55) node[above] {$\eta$};
    % C_infty: 竖直带状锥 {0} x R^{n-k}
    \fill[gray!16] (-0.42, -2.30) rectangle (0.42, 2.55);
    \node[left, inner sep=2pt] at (-0.46, 1.90) {\footnotesize $C_\infty$};
    % 两个锥 C_{eta_1}, C_{eta_2}
    \fill[gray!14] (0, 0) -- (4.20, 2.30) -- (4.20, 1.60) -- cycle;
    \fill[gray!14] (0, 0) -- (4.20, -1.60) -- (4.20, -2.30) -- cycle;
    \node[right, inner sep=2pt] at (4.24, 1.92) {\footnotesize $C_{\eta_1}$};
    \node[right, inner sep=2pt] at (4.24, -1.92) {\footnotesize $C_{\eta_2}$};
    % 直线 xi = xi_0 上位于 C_infty 之外的有界段
    \draw[densely dashed] (2.62, -2.30) -- (2.62, 2.55);
    \draw[<->, thick] (2.62, -0.94) -- (2.62, 0.94);
    \node[right, inner sep=2pt] at (2.68, 0.10) {\footnotesize 有界};
    \node[below, inner sep=1pt] at (2.62, -2.30) {\footnotesize $\xi = \xi_0$};
\end{tikzpicture}%
}

% 6. 图 6.6: 径向紧化中边界点的邻域基 U_delta
\newcommand{\drawFigSixSix}{%
\begin{tikzpicture}[scale=0.92, >=Stealth, line join=round]
    % 左图: 圆盘中的边界点 alpha 与 U_delta
    \begin{scope}
        \draw[thick] (0, 0) circle (2.05);
        \fill[gray!18] (2.05, 0) arc (0:42:2.05) arc (42:0:1.62) -- cycle;
        \fill (2.05, 0) circle (2.0pt) node[right, inner sep=1pt] {$\alpha$};
        \node at (1.35, 0.42) {\footnotesize $U_\delta$};
        \node[below] at (0, -2.20) {\footnotesize $\overline{\R^n}$};
    \end{scope}
    % 右图: U_delta 与 R^n 的交 (锥的远端)
    \begin{scope}[xshift=6.6cm]
        \draw[->] (-2.10, 0) -- (2.30, 0) node[right] {\footnotesize $x_1$};
        \draw[->] (0, -1.95) -- (0, 1.95);
        \fill[gray!18] (0, 0) -- (1.95, 1.32) -- (2.25, 1.32) -- (2.25, -1.32) -- (1.95, -1.32) -- cycle;
        \fill[gray!28] (1.55, 0) -- (2.25, 0.72) -- (2.25, -0.72) -- cycle;
        \node[right] at (2.28, 0.95) {\footnotesize $|x| > \delta^{-1}$};
        \node[below] at (0.30, -1.98) {\footnotesize $U_\delta \cap \R^n$};
    \end{scope}
\end{tikzpicture}%
}

% 7. 图 6.7: 径向紧化 R^n 与射影坐标的水平集 (n = 2)
\newcommand{\drawFigSixSeven}{%
\begin{tikzpicture}[scale=1.00, >=Stealth, line join=round]
    \draw[thick] (0, 0) circle (2.15);
    % 射影坐标 rho_1 = 1/x_1 的水平集: 竖直线
    \foreach \t in {0.45, 1.05, 1.65}{
        \pgfmathsetmacro{\h}{sqrt(2.15*2.15 - \t*\t)}
        \draw[gray!70] (\t, -\h) -- (\t, \h);
    }
    % 射影坐标 xhat_2 = x_2/x_1 的水平集: 过原点的射线
    \foreach \a in {30, 60, 120, 150}{
        \draw[gray!70] (0, 0) -- (\a:2.15);
    }
    \draw[->] (-2.36, 0) -- (2.36, 0);
    \draw[->] (0, -2.36) -- (0, 2.36);
    \node[above right, inner sep=1pt] at (1.62, 1.48) {\footnotesize $\rho_1 = c$};
    \node[above right, inner sep=1pt] at (0.62, 1.86) {\footnotesize $\hat{x}_2 = c$};
    \node[below] at (0, -2.42) {\footnotesize $\partial\overline{\R^2} = \mathbb{S}^{1}$};
    \fill (0, 0) circle (1.4pt) node[below left, inner sep=1pt] {$0$};
\end{tikzpicture}%
}

% 8. 图 6.8: 带角流形 [0, 1) x R^N
\newcommand{\drawFigSixEight}{%
\begin{tikzpicture}[scale=0.94, >=Stealth, line join=round]
    \fill[gray!10] (0, 0) rectangle (6.10, 2.30);
    \draw[thick] (0, 0) rectangle (6.10, 2.30);
    % 底边界 {0} x R^N
    \draw[line width=3.0pt] (0, 0) -- (6.10, 0);
    \node[below] at (3.05, -0.10) {\footnotesize $\{0\} \times \overline{\R^N}$ (半经典微局部对象所在)};
    % 右边界 [0,1) x dR^N
    \draw[line width=1.6pt, densely dashed] (6.10, 0) -- (6.10, 2.30);
    \node[right, align=left] at (6.16, 1.15) {\footnotesize $[0,1) \times \partial\overline{\R^N}$};
    % 坐标轴
    \draw[->] (0, 0) -- (-0.55, 0);
    \draw[->] (0, 0) -- (0, 2.72) node[above] {\footnotesize $h$};
    \node[below left, inner sep=1pt] at (-0.02, -0.02) {\footnotesize $0$};
    \node[above right, inner sep=1pt] at (0.12, 2.10) {\footnotesize $h \to 1$};
\end{tikzpicture}%
}

% 9. 图 6.9: 球极投影 (n = 1)
\newcommand{\drawFigSixNine}{%
\begin{tikzpicture}[scale=1.05, >=Stealth, line join=round]
    \draw[thick] (0, 0) circle (1.85);
    \fill[gray!12] (1.85, 0) arc (0:180:1.85) -- cycle;
    % 坐标轴
    \draw[->] (-2.45, 0) -- (3.10, 0) node[right] {\footnotesize $x_1$};
    \draw[->] (0, -0.45) -- (0, 2.35) node[above] {\footnotesize $x_2$};
    % 上半圆上的一点 x
    \coordinate (P) at (45:1.85);
    \draw[very thick] (0, 0) -- (P);
    \fill (P) circle (2.0pt);
    \node[right, inner sep=2pt] at (P) {\footnotesize $x$};
    % x_1, x_2 两条直角边
    \draw[densely dashed] (P) -- (1.308, 0);
    \draw[densely dashed] (P) -- (0, 1.308);
    \node[below, inner sep=2pt] at (0.90, 0) {\footnotesize $x_1$};
    \node[left, inner sep=2pt] at (0, 0.95) {\footnotesize $x_2$};
    % 赤道两端 (其像分别为 +infinity 与 -infinity)
    \fill (1.85, 0) circle (2.2pt);
    \node[right, inner sep=2pt] at (1.93, 0.10) {\footnotesize $+\infty$};
    \fill (-1.85, 0) circle (2.2pt);
    \node[left, inner sep=2pt] at (-1.93, 0.10) {\footnotesize $-\infty$};
    % 像: phi(x) = x_1/x_2
    \node[right, align=left, inner sep=2pt] at (2.05, 1.55) {\footnotesize $\varphi(x) = \dfrac{x_1}{x_2}$};
    \draw[->, gray!70] (2.02, 1.38) -- (1.55, 0.95);
    \node[below] at (0, -2.05) {\footnotesize $\mathbb{S}^1_{\ge}$};
\end{tikzpicture}%
}
